\documentclass[11pt,a4paper]{article}
\usepackage[utf8]{inputenc}
\usepackage[T1]{fontenc}
\usepackage{lmodern}
\usepackage[margin=2.2cm]{geometry}
\usepackage{amsmath,amssymb,amsthm,mathtools}
\usepackage{booktabs}
\usepackage{enumitem}
\usepackage[colorlinks=true,linkcolor=blue!60!black,citecolor=blue!60!black,urlcolor=blue!60!black]{hyperref}

\newtheorem{theorem}{Theorem}[section]
\newtheorem{corollary}[theorem]{Corollary}
\newtheorem{proposition}[theorem]{Proposition}
\theoremstyle{definition}
\newtheorem{definition}[theorem]{Definition}
\theoremstyle{remark}
\newtheorem{remark}[theorem]{Remark}

\title{\textbf{Collatz Trajectory Dynamics on Natural Modular Rings $\mathbb{Z}_{2^a}$ and $\mathbb{Z}_{2^a - 3^b}$:\\ $\mathrm{SL}_2(\mathbb{Z})$ Catalan Exponent-Basis Decomposition, Drift Envelope Bounds, and Sign-Free Positive Wronskian Sums in Lean~4}}
\author{Dmitry Karataev \and Luch (AI-Guiders)}
\date{2026}

\begin{document}
\maketitle

\begin{abstract}
We reformulate the Collatz $3n+1$ problem as a discrete trajectory dynamics problem on the two-sorted algebraic carrier structure $\langle\mathcal{M},\mathcal{T}\rangle$ (in the sense of A.\,I.~Mal'cev) over the natural modular rings $\mathbb{Z}_{2^a} = \mathbb{Z}_{4Q}$ ($Q = 2^{a-2}$) and the Pillai gap ring $\mathbb{Z}_{C_{a,b}}$ ($C_{a,b} = 2^a - 3^b$). First, extending the piecewise-linear rotator and Weierstrass stereographic calculus to an arbitrary four-quadrant ring $\mathbb{Z}_{4Q}$ (\texttt{thm18} in \texttt{RotatorCalculusPaper.lean}), we show that every $a$-step binary--ternary trajectory $2^a n_{\mathrm{final}} - 3^b n_{\mathrm{start}} = D_{a,b}$ decomposes into an exact quotient--remainder balance $2^a(n_{\mathrm{final}} - n_{\mathrm{start}}) + (2^a - 3^b)n_{\mathrm{start}} \equiv D_{a,b}$, where the drift polynomial $D_{a,b} = \sum_{j=0}^{b-1} 3^{b-1-j} 2^{s_j}$ is bounded unconditionally by the Pillai envelope $3^b - 2^b \le D_{a,b} \le 2^{a-b}(3^b - 2^b)$ and the $a$-bit residue $r = n_{\mathrm{start}}\bmod 2^a\in\mathbb{Z}_{2^a}$ bijectively determines the parity vector (contracting on $382/512 = 74.61\%$ of residues for $a=9$). Second, for any two consecutive trajectory arcs $P_1 n_{\mathrm{mid}} - Q_1 n_{\mathrm{start}} = D_1$ and $P_2 n_{\mathrm{final}} - Q_2 n_{\mathrm{mid}} = D_2$, both the composite drift $D_{\mathrm{tot}} = Q_2 D_1 + P_1 D_2$ and the composite Pillai gap $C_{\mathrm{tot}} = Q_2 C_1 + P_1 C_2$ are generated by the same linear operator $(Q_2, P_1)$, locking the chord displacement on any closed cycle ($n_{\mathrm{final}} = n_{\mathrm{start}}$) to the $2\times 2$ Wronskian determinant $C_{\mathrm{tot}}(n_{\mathrm{mid}} - n_{\mathrm{start}}) \equiv C_2 D_1 - C_1 D_2$. Third, we prove that the unit-determinant matrix $\mathcal{B}_{2,3} = \begin{pmatrix} 2 & 3 \\ 1 & 2 \end{pmatrix} \in \mathrm{SL}_2(\mathbb{Z})$ decomposes every exponent pair $(a,b)\in\mathbb{Z}^2$ uniquely into the Catalan basis $u = 2a - 3b$ (arc $(2,1)$ with gap $2^2 - 3^1 = +1$) and $v = 2b - a$ (Catalan arc $(3,2)$ with gap $3^2 - 2^3 = +1$), so that $2^a = 4^u 8^v$ and $3^b = 3^u 9^v$ with $u \ge 1, v \ge 1$ on every contracting cycle ($\frac{3}{2} < \frac{a}{b} < 2$). Splitting any closed cycle at the boundary between the $(4^u,3^u)$ block ($G_1 = 4^u - 3^u > 0$) and the $(8^v,9^v)$ block ($G_2 = 9^v - 8^v > 0$) transforms the Wronskian chord equation into a strictly positive sum without subtraction: $(2^a - 3^b)(n_{\mathrm{start}} - n_{\mathrm{mid}}) \equiv (9^v - 8^v)D_1 + (4^u - 3^u)D_2 > 0$ (\texttt{thm19}, \texttt{0 sorryAx}, \texttt{0 Classical.choice}).
\end{abstract}

\section{Introduction and Two-Sorted Carrier Formulation on $\mathbb{Z}_{2^a}$ and $\mathbb{Z}_{2^a - 3^b}$}
\label{sec:intro}

Throughout this paper, we follow A.\,I.~Mal'cev~\cite{Malcev1973} in defining a \emph{carrier} (\emph{носитель}) as the underlying domain set of a sort in a two-sorted algebraic structure $\langle\mathcal{M},\mathcal{T}\rangle$. In our foundational rotator calculus~\cite{KarataevRotator1}, the discrete sort $\mathcal{M}$ was instantiated at the self-closed register modulus $2^9 = 512$ ($\mathbb{Z}_{512}$). Crucially, the piecewise-linear rotator $(\mathcal{C}_{4Q},\mathcal{S}_{4Q})$, its constant difference Wronskian $\mathcal{W}_{4Q} \equiv Q$, and the Weierstrass stereographic embedding hold without modification on \textbf{any} four-quadrant modular ring $\mathbb{Z}_{4Q}$ ($Q\in\mathbb{Z}$, Theorem~\ref{thm:arbitrary-modulus-wronskian}).

Consider the compressed Collatz operator on positive integers:
\begin{equation}
  \label{eq:collatz-step}
  T(x) \;=\;
  \begin{cases}
    x / 2, & x \equiv 0 \pmod 2, \\
    (3x + 1) / 2, & x \equiv 1 \pmod 2.
  \end{cases}
\end{equation}
Over $a \ge 1$ binary divisions involving $b \in \{0,\dots,a\}$ odd steps at indices $0 \le s_0 < s_1 < \dots < s_{b-1} \le a - 1$, the trajectory from $n_{\mathrm{start}}$ to $n_{\mathrm{final}} = T^{(a)}(n_{\mathrm{start}})$ satisfies the exact linear Diophantine equation of motion:
\begin{equation}
  \label{eq:collatz-equation-of-motion}
  2^a \cdot n_{\mathrm{final}} \;-\; 3^b \cdot n_{\mathrm{start}} \;=\; D_{a,b}, \qquad
  D_{a,b} \;:=\; \sum_{j=0}^{b-1} 3^{b-1-j}\,2^{s_j}.
\end{equation}
Writing $n_{\mathrm{start}} = 2^a q + r$ with quotient $q = \lfloor n_{\mathrm{start}}/2^a\rfloor \ge 0$ and phase residue $r = n_{\mathrm{start}}\bmod 2^a \in \mathbb{Z}_{2^a}$, the parity sequence $(e_0,\dots,e_{a-1})\in\{0,1\}^a$ is a bijection of $r\in\mathbb{Z}_{2^a}$ independent of $q$. Consequently, $b = b(r)$ follows the binomial Pascal distribution $\binom{a}{b}$ across $\mathbb{Z}_{2^a}$ (for $a=9$ on $\mathbb{Z}_{512}$, $382/512 = 74.61\%$ of residues have $b(r)\le 5$ and $3^{b(r)}\le 243 < 512$), and
\begin{equation}
  \label{eq:block-affine-transition}
  T^{(a)}(2^a q + r) \;=\; 3^{b(r)}\,q \;+\; x_a(r), \qquad
  2^a\,x_a(r) \;=\; 3^{b(r)}\,r \;+\; D_{a,b}(r).
\end{equation}
Moreover, across all $\binom{a}{b}$ subsets $0 \le s_0 < \dots < s_{b-1} \le a-1$, the drift polynomial $D_{a,b}$ is tightly bracketed by the geometric Pillai envelope:
\begin{equation}
  \label{eq:drift-envelope}
  3^b - 2^b \;=\; \sum_{j=0}^{b-1} 3^{b-1-j}\,2^j \;\le\; D_{a,b} \;\le\; 2^{a-b}\sum_{j=0}^{b-1} 3^{b-1-j}\,2^j \;=\; 2^{a-b}\,\bigl(3^b - 2^b\bigr).
\end{equation}

\section{Arbitrary-Modulus Rotator Calculus on $\mathbb{Z}_{4Q}$ and the Two-Arc Wronskian Law}
\label{sec:two-arc-wronskian}

\begin{theorem}[Arbitrary-modulus rotator embedding on $\mathbb{Z}_{4Q}$ and two-arc trajectory Wronskian law, \texttt{thm18} in \texttt{RotatorCalculusPaper.lean}]
\label{thm:arbitrary-modulus-wronskian}
For all integers $Q, A, B, C, u, v, w, \varphi_0, P_1, Q_1, P_2, Q_2, n_{\mathrm{start}}, n_{\mathrm{mid}}, n_{\mathrm{final}}, D_1, D_2 \in \mathbb{Z}$:
\begin{enumerate}[label=\textbf{(\roman*)},leftmargin=*]
\item \textbf{Arbitrary-modulus rotator and stereographic invariant on $\mathbb{Z}_{4Q}$:}
  for $\mathcal{C}_{4Q}(\varphi_0) = Q - \varphi_0$ and $\mathcal{S}_{4Q}(\varphi_0) = \varphi_0$, the $L^1$ norm and difference Wronskian satisfy $\mathcal{C}_{4Q}(\varphi_0) + \mathcal{S}_{4Q}(\varphi_0) \equiv Q$ and $\mathcal{W}_{4Q}(\varphi_0) = \mathcal{C}_{4Q}(\varphi_0)\mathcal{S}_{4Q}(\varphi_0+1) - \mathcal{S}_{4Q}(\varphi_0)\mathcal{C}_{4Q}(\varphi_0+1) \equiv Q$. For any triple $A + B = C$ with Bézout partner $A v - B u = 1$ ($u + v = w$), setting $I_{\mathrm{rot},Q} = B\mathcal{C}_{4Q}(\varphi_0) - A\mathcal{S}_{4Q}(\varphi_0)$ and $\Delta_{\mathrm{rot},Q} = \mathcal{C}_{4Q}(\varphi_0)v - \mathcal{S}_{4Q}(\varphi_0)u$ yields
  \begin{equation}
    \label{eq:arbitrary-q-invariant}
    Q\,A \;=\; C\,\mathcal{C}_{4Q}(\varphi_0) - I_{\mathrm{rot},Q}, \qquad
    Q\,B \;=\; C\,\mathcal{S}_{4Q}(\varphi_0) + I_{\mathrm{rot},Q}, \qquad
    C\cdot\Delta_{\mathrm{rot},Q} \;-\; w\cdot I_{\mathrm{rot},Q} \;\equiv\; Q;
  \end{equation}
\item \textbf{Two-arc composition and displacement decomposition in $\mathbb{Z}_{C_{\mathrm{tot}}}$:}
  given two consecutive trajectory arcs $P_1 n_{\mathrm{mid}} - Q_1 n_{\mathrm{start}} = D_1$ and $P_2 n_{\mathrm{final}} - Q_2 n_{\mathrm{mid}} = D_2$ (with $P_i = 2^{a_i},\ Q_i = 3^{b_i},\ C_i = P_i - Q_i,\ C_{\mathrm{tot}} = P_1 P_2 - Q_1 Q_2,\ D_{\mathrm{tot}} = Q_2 D_1 + P_1 D_2$), the composite trajectory satisfies
  \begin{equation}
    \label{eq:composite-displacement}
    (P_1 P_2)\,n_{\mathrm{final}} \;-\; (Q_1 Q_2)\,n_{\mathrm{start}} \;=\; D_{\mathrm{tot}}, \qquad
    (P_1 P_2)\,\bigl(n_{\mathrm{final}} - n_{\mathrm{start}}\bigr) \;+\; C_{\mathrm{tot}}\,n_{\mathrm{start}} \;\equiv\; D_{\mathrm{tot}};
  \end{equation}
\item \textbf{Cramer cycle lock and $2\times 2$ Wronskian chord identity:}
  on any closed cycle ($n_{\mathrm{final}} = n_{\mathrm{start}}$), both vertices $n_{\mathrm{start}}$ and $n_{\mathrm{mid}}$ are locked over the Pillai denominator $C_{\mathrm{tot}}$, and their chord displacement equals the $2\times 2$ Wronskian determinant between the arc Pillai gaps and arc drifts:
  \begin{equation}
    \label{eq:cycle-wronskian-chord}
    C_{\mathrm{tot}}\,n_{\mathrm{start}} = Q_2 D_1 + P_1 D_2, \qquad
    C_{\mathrm{tot}}\,n_{\mathrm{mid}} = P_2 D_1 + Q_1 D_2, \qquad
    C_{\mathrm{tot}}\,\bigl(n_{\mathrm{mid}} - n_{\mathrm{start}}\bigr) \;\equiv\; C_2 D_1 - C_1 D_2.
  \end{equation}
\end{enumerate}
\end{theorem}

\section{$\mathrm{SL}_2(\mathbb{Z})$ Catalan Exponent-Basis Decomposition and Sign-Free Positive Wronskian Sum}
\label{sec:catalan-basis}

In general two-arc splittings~\eqref{eq:cycle-wronskian-chord}, if $C_1$ and $C_2$ have the same sign, the right-hand side $C_2 D_1 - C_1 D_2$ involves a subtraction. We now eliminate this sign ambiguity completely by decomposing the exponent vector $(a,b)$ in the canonical $\mathrm{SL}_2(\mathbb{Z})$ basis generated by the two unit-gap Catalan locks $E_0 = (2,1)$ ($2^2 - 3^1 = +1$) and $E_1 = (3,2)$ ($3^2 - 2^3 = +1$).

\begin{theorem}[$\mathrm{SL}_2(\mathbb{Z})$ Catalan exponent-basis decomposition and sign-free positive Wronskian sum on closed cycles, \texttt{thm19} in \texttt{RotatorCalculusPaper.lean}]
\label{thm:catalan-basis-positive-wronskian}
For any exponent pair $(a,b)\in\mathbb{Z}^2$ and any two consecutive trajectory arcs $P_1 n_{\mathrm{mid}} - Q_1 n_{\mathrm{start}} = D_1$ and $P_2 n_{\mathrm{final}} - Q_2 n_{\mathrm{mid}} = D_2$:
\begin{enumerate}[label=\textbf{(\roman*)},leftmargin=*]
\item \textbf{Unimodular Catalan basis $(2,1)$ and $(3,2)$ in $\mathrm{SL}_2(\mathbb{Z})$:}
  the basis matrix $\mathcal{B}_{2,3} = \begin{pmatrix} 2 & 3 \\ 1 & 2 \end{pmatrix}$ has determinant $\det\mathcal{B}_{2,3} = 2\times 2 - 3\times 1 = +1$, uniquely decomposing every $(a,b)\in\mathbb{Z}^2$ into
  \begin{equation}
    \label{eq:catalan-uv-coords}
    u \;:=\; 2a - 3b, \qquad v \;:=\; 2b - a, \qquad
    \begin{pmatrix} a \\ b \end{pmatrix} \;=\; u\begin{pmatrix} 2 \\ 1 \end{pmatrix} + v\begin{pmatrix} 3 \\ 2 \end{pmatrix}, \qquad
    u + v \;=\; a - b.
  \end{equation}
  For every contracting Collatz cycle ($2^a > 3^b \implies \frac{a}{b} > \log_2 3 > \frac{3}{2}$, and $a < 2b \implies \frac{a}{b} < 2$), both coordinates are strictly positive integers ($u \ge 1,\ v \ge 1$), factoring the total powers without fractions into
  \begin{equation}
    \label{eq:power-factorization-4-3-8-9}
    2^a \;=\; (2^2)^u\cdot(2^3)^v \;=\; 4^u\cdot 8^v, \qquad
    3^b \;=\; (3^1)^u\cdot(3^2)^v \;=\; 3^u\cdot 9^v;
  \end{equation}
\item \textbf{Decomposition of the Pillai gap $C_{\mathrm{tot}} = 2^a - 3^b$ into unit-gap differences $G_1 = 4^u - 3^u > 0$ and $G_2 = 9^v - 8^v > 0$:}
  setting $(P_1,Q_1) = (4^u,3^u)$ ($G_1 := P_1 - Q_1 = 4^u - 3^u > 0$, divisible by $4 - 3 = 1$) and $(P_2,Q_2) = (8^v,9^v)$ ($G_2 := Q_2 - P_2 = 9^v - 8^v > 0$, divisible by $9 - 8 = 1$), the total Pillai gap $C_{\mathrm{tot}} = P_1 P_2 - Q_1 Q_2$ satisfies
  \begin{equation}
    \label{eq:gap-split-g1-g2}
    C_{\mathrm{tot}} \;\equiv\; 8^v\,\bigl(4^u - 3^u\bigr) \;-\; 3^u\,\bigl(9^v - 8^v\bigr) \;\equiv\; 9^v\,\bigl(4^u - 3^u\bigr) \;-\; 4^u\,\bigl(9^v - 8^v\bigr);
  \end{equation}
\item \textbf{Sign-free strictly positive Wronskian chord sum on any closed cycle ($n_{\mathrm{final}} = n_{\mathrm{start}}$):}
  splitting any closed cycle at the intermediate vertex $n_{\mathrm{mid}}$ after the first $2u$ binary steps ($b_1 = u$ odd steps, $P_1 = 4^u,\ Q_1 = 3^u$, drift $D_1 \ge 1$; remaining arc $a_2 = 3v,\ b_2 = 2v$, $P_2 = 8^v,\ Q_2 = 9^v$, drift $D_2 \ge 1$) gives opposite signs to $C_1 = +G_1 > 0$ and $C_2 = -G_2 < 0$, converting the Wronskian identity~\eqref{eq:cycle-wronskian-chord} into a \textbf{sum of two strictly positive terms}:
  \begin{equation}
    \label{eq:positive-wronskian-sum}
    \bigl(2^a - 3^b\bigr)\cdot\bigl(n_{\mathrm{start}} - n_{\mathrm{mid}}\bigr) \;\equiv\; \underbrace{\bigl(9^v - 8^v\bigr)\,D_1}_{> 0} \;+\; \underbrace{\bigl(4^u - 3^u\bigr)\,D_2}_{> 0} \;>\; 0,
  \end{equation}
  which unconditionally forces $n_{\mathrm{start}} > n_{\mathrm{mid}}$ and bounds both $(9^v - 8^v)D_1$ and $(4^u - 3^u)D_2$ from above by $(2^a - 3^b)(n_{\mathrm{start}} - n_{\mathrm{mid}})$.
\end{enumerate}
\end{theorem}

\begin{table}[htbp]
\centering
\caption{First six resonant Stern--Brocot exponent stages of $\log_2 3$ and their Catalan $(u,v)$ coordinates.}
\label{tab:collatz-stages}
\small
\begin{tabular}{@{}llllll@{}}
\toprule
\textbf{Stage} & \textbf{Exponents $(a,b)$} & \textbf{Catalan $(u,v)$} & \textbf{Powers $(2^a, 3^b)$} & \textbf{Pillai Gap $C_{a,b} = 2^a - 3^b$} & \textbf{Linear Combination} \\
\midrule
$E_0$ & $(2, 1)$ & $(1, 0)$ & $(4, 3)$ & $+1$ & Trivial cycle $n = 1/(4-3) = 1$ \\
$E_1$ & $(3, 2)$ & $(0, 1)$ & $(8, 9)$ & $-1$ & Catalan Lock~I ($3^2 - 2^3 = 1$) \\
$E_2$ & $(5, 3)$ & $(1, 1)$ & $(32, 27)$ & $+5$ & $3^2(+1) + 2^2(-1) = 9 - 4 = +5$ \\
$E_3$ & $(8, 5)$ & $(1, 2)$ & $(256, 243)$ & $+13$ & $3^3(-1) + 2^3(+5) = -27 + 40 = +13$ \\
$E_4$ & $(11, 7)$ & $(1, 3)$ & $(2048, 2187)$ & $-139$ & $3^5(-1) + 2^3(+13) = -243 + 104 = -139$ \\
$E_5$ & $(19, 12)$ & $(2, 5)$ & $(524288, 531441)$ & $-7153$ & $3^7(+13) + 2^8(-139) = -7153$ \\
\bottomrule
\end{tabular}
\end{table}

\section{Formal Verification in Lean~4}
Theorems~\ref{thm:arbitrary-modulus-wronskian} and~\ref{thm:catalan-basis-positive-wronskian} (\texttt{thm18} and \texttt{thm19}) are formally verified in \texttt{RotatorCalculusPaper.lean} at \url{https://github.com/AI-Guiders/scientific-papers} (\texttt{0 sorryAx}, \texttt{0 Classical.choice}).

\begin{thebibliography}{9}
\bibitem{Malcev1973}
A.\,I.~Mal'cev, \emph{Algebraic Systems}, Die Grundlehren der mathematischen Wissenschaften, vol.~192, Springer-Verlag, 1973 (Russian original: Nauka, Moscow, 1970).
\bibitem{KarataevRotator1}
D.~Karataev and Luch, ``Cuboctahedral Rotator Calculus: Discrete Euler Exponential on $\mathbb{Z}_{512}$, Universal Constant Wronskian, and Closed-Form Quadrature on the Cuboctahedron $O_2(3)$ in Lean~4,'' \url{https://github.com/AI-Guiders/scientific-papers}, 2026.
\bibitem{Lagarias2010}
J.\,C.~Lagarias (ed.), \emph{The Ultimate Challenge: The $3x+1$ Problem}, American Mathematical Society, Providence, RI, 2010.
\bibitem{Moura2021}
L.~de~Moura and S.~Ullrich, ``The Lean~4 theorem prover and programming language,'' in \emph{CADE~28}, LNCS vol.~12699, Springer, 2021, pp.~625--635.
\end{thebibliography}

\end{document}
